\documentclass[11pt]{article}

% --- Packages for math & formatting ---
\usepackage{amsmath, amssymb, amsthm} % advanced math
\usepackage{listings}                  % code formatting
\usepackage{xcolor}                     % color for code
\usepackage{graphicx}                   % include figures
\usepackage{hyperref}                   % clickable links
\usepackage{authblk}                    % author affiliations
\usepackage[margin=1in]{geometry}
\usepackage{fvextra} % better verbatim with line wrapping
\usepackage{float}

% --- Configure code display ---
\lstset{
  language=R,
  basicstyle=\ttfamily\small,
  numbers=left,
  stepnumber=1,
  numbersep=5pt,
  frame=single,
  backgroundcolor=\color{gray!10},
  breaklines=true,
  captionpos=b
}

% --- Theorem environment ---
\newtheorem{theorem}{Theorem}
\newtheorem{lemma}{Lemma}

% --- New commands ---
\newcommand*\mean[1]{\overline{#1}}
\newcommand{\var}{\text{var}}
\newcommand{\cov}{\text{cov}}
\newcommand{\cor}{\text{cor}}
\newcommand{\Rp}{\text{Re}}
\newcommand{\E}{\text{E}}
\newcommand{\ltsgr}{\text{ltsgr}}
\newcommand{\expit}{\text{expit}}
\newcommand{\logit}{\text{logit}}

% --- Header material ---
\title{How to fit models with species occurrence data using \texttt{xsdm}}
\author[1]{Daniel C. Reuman}
\author[1]{Angel Lu\'{i}s Robles Fern\'{a}ndez}
\affil[1]{Department of Ecology and Evolutionary Biology and Center for Ecological Research,
University of Kansas}
\date{}

% --- Make Sweave code/output blocks obey text margins and wrap nicely ---
\DefineVerbatimEnvironment{Sinput}{Verbatim}{
  fontsize=\small,
  xleftmargin=0pt,
  xrightmargin=0pt,
  breaklines=true,
  breakanywhere=true,
  commandchars=\\\{\}
}

\DefineVerbatimEnvironment{Soutput}{Verbatim}{
  fontsize=\small,
  xleftmargin=0pt,
  xrightmargin=0pt,
  breaklines=true,
  breakanywhere=true,
  commandchars=\\\{\}
}

% Optional: if you have Scode blocks too
\DefineVerbatimEnvironment{Scode}{Verbatim}{
  fontsize=\small,
  xleftmargin=0pt,
  xrightmargin=0pt,
  breaklines=true,
  breakanywhere=true,
  commandchars=\\\{\}
}

% --- start the doc ---
\usepackage{Sweave}
\begin{document}
\input{howToFit-knitr-concordance}

\begin{Schunk}
\begin{Sinput}
> library(knitr)
> opts_chunk$set(
+ concordance=TRUE
+ )
\end{Sinput}
\end{Schunk}


\begin{Schunk}
\begin{Sinput}
> knitr::opts_chunk$set(
+   echo=TRUE,
+   message=FALSE,
+   warning=FALSE,
+   cache=FALSE
+ )
\end{Sinput}
\end{Schunk}

\maketitle

\begin{abstract}
\noindent After reading the document entitled "The xsdm model", this document
introduces the statistically and computationally sophisticated user to the 
main functions of the \texttt{xsdm} package. The package implements a frequentist 
analysis of the xsdm model. This document should get the user to the point
of maximizing the likelihood of the model and profiling in ``easy'' cases, 
with callouts to
several troubleshooting documents about what to do when the process does
not go smoothly. Alternative fitted models can also be compared by AIC or another
criterion. The example worked out here is for a virtual species (i.e., generated 
data), so we also demonstrate that fitting an xsdm model can recover the true
parameters if model mis-specification is not an issue.
\end{abstract}

\section{Introduction to the running example used throughout this document}

This manual uses a small built-in example included with the \texttt{xsdm} R 
package to illustrate a baseline workflow. The example contains (i) an
occurrence table for a virtual species, and (ii) two environmental
\texttt{terra} raster time series representing CHELSA-derived 
bioclimatic variables for a region of interest over 39 years.

Load the data and get oriented to the first bioclimatic variable, which is mean annual temperature
in units of $100 \times ^\circ$C for 39 years for a particular region of interest. 
And make it units of $^\circ$C: 
\begin{Schunk}
\begin{Sinput}
> library(xsdm)
> library(terra)
> data("example_1")
> bio_1 <- terra::unwrap(example_1$bio01) 
> class(bio_1)
\end{Sinput}
\begin{Soutput}
[1] "SpatRaster"
attr(,"package")
[1] "terra"
\end{Soutput}
\begin{Sinput}
> dim(bio_1)
\end{Sinput}
\begin{Soutput}
[1] 128 123  39
\end{Soutput}
\begin{Sinput}
> mm <- terra::global(bio_1,fun=c("min","max"))
> min(mm$min) #global min over all years and all locations
\end{Sinput}
\begin{Soutput}
[1] -154
\end{Soutput}
\begin{Sinput}
> max(mm$max) #global max
\end{Sinput}
\begin{Soutput}
[1] 2061
\end{Soutput}
\begin{Sinput}
> bio_1 <- bio_1/100 #Make it degrees C
\end{Sinput}
\end{Schunk}

Now load the second environmental variable, which is annual total precipitation, in units of kg/m$^2$, for
the same 39 years and the same region of interest:
\begin{Schunk}
\begin{Sinput}
> bio_12 <- terra::unwrap(example_1$bio12) 
> class(bio_12)
\end{Sinput}
\begin{Soutput}
[1] "SpatRaster"
attr(,"package")
[1] "terra"
\end{Soutput}
\begin{Sinput}
> dim(bio_12)
\end{Sinput}
\begin{Soutput}
[1] 128 123  39
\end{Soutput}
\begin{Sinput}
> mm <-  terra::global(bio_12,fun=c("min","max"))
> min(mm$min) #again the global min, all years and locations
\end{Sinput}
\begin{Soutput}
[1] 45
\end{Soutput}
\begin{Sinput}
> max(mm$max) #global max
\end{Sinput}
\begin{Soutput}
[1] 1761
\end{Soutput}
\end{Schunk}

If the scales of your environmental variables are vastly different, it can cause problems
with optimization (see the document ``Troubleshooting: Optimization problems related to scaling").
So we change the units for precipitation to make the range of values similar to those for temperature.
Now the units are going to be dg/m$^2$:
\begin{Schunk}
\begin{Sinput}
> bio_12 = bio_12/100
> mm <- terra::global(bio_1,fun=c("min","max"))
> min(mm$min)
\end{Sinput}
\begin{Soutput}
[1] -1.54
\end{Soutput}
\begin{Sinput}
> max(mm$max)
\end{Sinput}
\begin{Soutput}
[1] 20.61
\end{Soutput}
\begin{Sinput}
> mm <- terra::global(bio_12,fun=c("min","max"))
> min(mm$min)
\end{Sinput}
\begin{Soutput}
[1] 0.45
\end{Soutput}
\begin{Sinput}
> max(mm$max)
\end{Sinput}
\begin{Soutput}
[1] 17.61
\end{Soutput}
\end{Schunk}

Next load data on detections and non-detections/pseudo-absences for a species:
\begin{Schunk}
\begin{Sinput}
> d <- example_1$occ_df
> class(d)
\end{Sinput}
\begin{Soutput}
[1] "tbl_df"     "tbl"        "data.frame"
\end{Soutput}
\begin{Sinput}
> dim(d)
\end{Sinput}
\begin{Soutput}
[1] 4000    4
\end{Soutput}
\begin{Sinput}
> head(d)
\end{Sinput}
\begin{Soutput}
               name      lon     lat presence
1 Species virtualis -1108723 1402222        1
2 Species virtualis  -883723 1402222        0
3 Species virtualis  -808723 1422222        1
4 Species virtualis -1168723 1552222        0
5 Species virtualis -1028723 1437222        0
6 Species virtualis  -948723 1452222        1
\end{Soutput}
\begin{Sinput}
> sum(d$presence == 1)
\end{Sinput}
\begin{Soutput}
[1] 1609
\end{Soutput}
\end{Schunk}
%DAN: Note, this chunk will change when the data for the species become wrapped in the package 

Now take a basic look at the environmental variables and the species data, just to see what we
have. Start
by getting the means through time of the environmental time series and plotting:
\begin{Schunk}
\begin{Sinput}
> m_bio_1 = terra::app(bio_1,mean)
> m_bio_12 = terra::app(bio_12,mean)
> par(mfrow=c(1,2))
> terra::plot(m_bio_1,axes=TRUE,main="Mean, bio1",xlab="x",ylab="y",legend=TRUE)
> terra::plot(m_bio_12,axes=TRUE,main="Mean, bio12",xlab="x",ylab="y",legend=TRUE)
\end{Sinput}
\end{Schunk}
%DAN: The spacing here is bad (way too much whitespace), and the sizing of the figures is sub-optimal
%They should be bigger. Not sure how to fix.

\noindent Now do the same for standard deviations:
\begin{Schunk}
\begin{Sinput}
> sd_bio_1 = terra::app(bio_1,sd)
> sd_bio_12 = terra::app(bio_12,sd)
> par(mfrow=c(1,2))
> terra::plot(sd_bio_1,axes=TRUE,main="SD, bio1",xlab="x",ylab="y",legend=TRUE)
> terra::plot(sd_bio_12,axes=TRUE,main="SD, bio12",xlab="x",ylab="y",legend=TRUE)
\end{Sinput}
\end{Schunk}

\noindent Now plot the species detections and non-detections/pseudo-absences on a backdrop of the mean of bio 1:
\begin{Schunk}
\begin{Sinput}
> d <- as.data.frame(d)
> pts_0 <- terra::vect(d[d$presence == 0, ],
+                      geom=c("lon","lat"),
+                      crs=terra::crs(m_bio_1))
> pts_1 <- terra::vect(d[d$presence == 1, ],
+                      geom = c("lon","lat"),
+                      crs=terra::crs(m_bio_1))
> par(mfrow=c(1,2))
> terra::plot(m_bio_1,
+             axes = TRUE,
+             main = "Non-detections",
+             xlab="x",
+             ylab="y",
+             legend = TRUE)
> terra::plot(pts_0, add=TRUE,col="black",pch=20,cex=.2)
> terra::plot(m_bio_1,
+             axes = TRUE,
+             main = "Detections",
+             xlab = "x",
+             ylab = "y",
+             legend = TRUE)
> terra::plot(pts_1, add=TRUE,col="black",pch=20,cex=.2)
\end{Sinput}
\end{Schunk}

The fitting tools offered by \texttt{xsdm} use a more succinct version of the data,
since only the environmental time series at the locations for which there is a 
detection or a non-detection/pseudo-absence matter; though we will return to using 
rasters after fitting and model selection. So, for now, we cut out only those
time series at the species locations and store them in an array using the convenience 
function \texttt{env\_data\_array}:

\begin{Schunk}
\begin{Sinput}
> env_data <- list(bio_1 = bio_1, bio_12 = bio_12)  # order defines variable 1 and 2
> env_array <- env_data_array(env_data, d)  # (n locations) x (time) x (p vars)
> class(env_array)
\end{Sinput}
\begin{Soutput}
[1] "array"
\end{Soutput}
\begin{Sinput}
> dim(env_array) 
\end{Sinput}
\begin{Soutput}
[1] 4000   39    2
\end{Soutput}
\end{Schunk}

\noindent And we pull out the presences/pseudo-absences as a binary vector:
\begin{Schunk}
\begin{Sinput}
> occ <- d$presence
> length(occ)
\end{Sinput}
\begin{Soutput}
[1] 4000
\end{Soutput}
\end{Schunk}

\noindent Now we are ready to think about how the xsdm model applies to our data.

\section{Evaluating the likelihood}

The xsdm model's likelihood function is described in the document "The xsdm model",
where model parameter are also described in detail. Parameters are:
\begin{itemize}
\item $\vec{\mu}$, which encodes ideal values for population growth 
of the two environmental variables (a length-2 vector for the present 
scenario of two environmental variables, unconstrained values); 
\item $\tilde{\vec{\sigma}}_L$, which has to do with breadth of the growth-environment
function ``to the left'' with respect to two basis vectors (a length-2 vector for our case,
positive entries);
%Note, here I am NOT saying these can be Inf, will cover the boundary models in another doc
\item $\tilde{\vec{\sigma}}_R$, which is similar but ``to the right'' (a length-2 vector,
positive entries);
\item $\tilde{c}$ and $p_d$, which relate to species detection (scalars, $\tilde{c}$ is 
unconstrained and $0<p_d\leq 1$); and
\item $O$, which is an orthogonal matrix which 
describes the basis vectors mentioned above ($2 \times 2$ for the present example, must be
orthogonal, i.e., $OO^{\tau}=I$ for $\tau$ the transpose).
\end{itemize}
In code, the above parameters are denoted \texttt{mu}, \texttt{sigltil}, \texttt{sigrtil},
\texttt{ctil}, \texttt{pd}, and \texttt{o\_mat}.

The likelihood is coded as \texttt{loglik\_bio}, so as an introduction to the function 
let's evaluate it at a haphazardly chosen set of parameters:
\begin{Schunk}
\begin{Sinput}
> mu <- c(1, 1)
> sigltil <- c(1, 4)
> sigrtil <- c(2, 3)
> ctil <- -10
> pd <- 0.8
> o_mat <- diag(2)
> loglik_bio(env_dat = env_array,
+            occ = occ,
+            mu = mu,
+            sigltil = sigltil,
+            sigrtil = sigrtil,
+            o_mat = o_mat,
+            ctil = ctil,
+            pd = pd)
\end{Sinput}
\begin{Soutput}
[1] -2110.297
\end{Soutput}
\end{Schunk}
\noindent By default, the function gives the log likelihood, but you can also get the linear-scale
likelihood:
\begin{Schunk}
\begin{Sinput}
> loglik_bio(env_array,
+            occ,
+            mu,
+            sigltil,
+            sigrtil,
+            o_mat,
+            ctil,
+            pd,
+            return_prob = TRUE)
\end{Sinput}
\begin{Soutput}
[1] 0
\end{Soutput}
\end{Schunk}
\noindent For these haphazardly chosen parameters, the (linear-scale) likelihood is zero to within 
numeric precision - this is typical. We next optimize.
 
\section{An unconstrained parameter space}
 
Some model parameters are constrained, i.e., only certain values are allowed (specifically, 
\texttt{sigltil}, \texttt{sigrtil}, \texttt{pd}, and \texttt{o\_mat}; see previous section). 
Rather than attempt to perform optimizations subject to these constraints, we 
here define a transformation from an unconstrained Euclidean space to 
the space of allowed parameters, and we subsequently optimize on the unconstrained space. 
Parameters in the unconstrained space are henceforth called
``math-scale'' parameters, and the constrained parameters described above are called 
``biological-scale'' parameters because they more directly represent biological quantities. 
When a distinction is needed in mathematical notation, we denote biological-scale parameters with a superscript
$(b)$, and math-scale parameters with a superscript $(m)$, i.e., $\vec{\mu}^{(b)}$, $O^{(b)}$, 
$\tilde{\vec{\sigma}}_L^{(b)}$,
$\tilde{\vec{\sigma}}_R^{(b)}$, $\tilde{c}^{(b)}$, and $p_d^{(b)}$ versus $\vec{\mu}^{(m)}$, $O^{(m)}$, 
$\tilde{\vec{\sigma}}_L^{(m)}$,
$\tilde{\vec{\sigma}}_R^{(m)}$, $\tilde{c}^{(m)}$, and $p_d^{(m)}$.

The transformation from math- to biological-scale parameters acts separately on each parameter, 
as follows for all the parameters except $O$:
\begin{align}
\vec{\mu}^{(b)} &= \vec{\mu}^{(m)} \label{eq:mutrans}\\
\tilde{\vec{\sigma}}_L^{(b)} &= \exp(\tilde{\vec{\sigma}}_L^{(m)}) \label{eq:sigLtrans}\\
\tilde{\vec{\sigma}}_R^{(b)} &= \exp(\tilde{\vec{\sigma}}_R^{(m)}) \label{eq:sigRtrans}\\
\tilde{c}^{(b)} &= \tilde{c}^{(m)} \label{eq:ctrans}\\
p_d^{(b)} &= \expit(p_d^{(m)}). \label{eq:pdtrans}
\end{align}
Here, $\exp$ of a vector is interpreted to be the vector resulting from
$\exp$-transforming each component, and $\expit$ is the standard logistic sigmoid
function $1/(1+\exp(-x))$, i.e., the inverse of the $\logit$ function.

The parameter $O^{(m)}$ is interpreted to be an unconstrained Euclidean vector
of dimension $\frac{p^2-p}{2}$, where $p$ is the number of environmental variables being considered
($p=2$ in our running example), and $O^{(b)}$ is obtained from $O^{(m)}$
by using $O^{(m)}$ to form a skew-symmetric matrix of dimensions $p \times p$, 
and then applying the matrix exponential map to that skew-symmetric matrix. It is known that
the matrix exponential of a skew-symmetric matrix is an orthogonal matrix.

The unconstrained space of parameters (the space of math-scale parameters) 
is thus a Euclidean space of dimensions $3p+\frac{p^2-p}{2}+2$. 
The $3p$ term in this expression comes from the parameters $\vec{\mu}^{(m)}$,
$\tilde{\vec{\sigma}}_L^{(m)}$, and $\tilde{\vec{\sigma}}_R^{(m)}$;
the $2$ in this expression comes from the parameters $\tilde{c}^{(m)}$ and
$p_d^{(m)}$; and the term $\frac{p^2-p}{2}$ in the expression comes from
$O^{(m)}$.

The transformation from math-scale to bio-scale parameters is implemented
in \texttt{xsdm} using the function \texttt{math\_to\_bio}, which takes an unconstrained
numeric vector as its argument and returns a named list of of biological-scale
parameters. But it is not so common for the end-user to call \texttt{math\_to\_bio}
directly because we have written the likelihood directly in terms of the math-scale
parameters in a function \texttt{loglik\_math}. We now demonstrate \texttt{math\_to\_bio}
and \texttt{loglik\_math}:
\begin{Schunk}
\begin{Sinput}
> #loglik_bio and math_to_bio require a named argument to help prevent errors - 
> #see below - make_mask_names helps set it up
> param_vector <- make_mask_names(p = 2) #p = 2 environmental variables in our case
> param_vector 
\end{Sinput}
\begin{Soutput}
     mu1      mu2 sigltil1 sigltil2 sigrtil1 sigrtil2     ctil       pd 
      NA       NA       NA       NA       NA       NA       NA       NA 
  o_par1 
      NA 
\end{Soutput}
\begin{Sinput}
> set.seed(101)
> param_vector[1:9] <- rnorm(9) #fill with random values just to try it
> math_to_bio(param_vector)
\end{Sinput}
\begin{Soutput}
$mu
[1] -0.3260365  0.5524619

$sigltil
[1] 0.509185 1.239068

$sigrtil
[1] 1.364474 3.234797

$ctil
[1] 0.6187899

$pd
[1] 0.4718462

$o_mat
          [,1]       [,2]
[1,] 0.6081818 -0.7937978
[2,] 0.7937978  0.6081818
\end{Soutput}
\begin{Sinput}
> loglik_math(param_vector, env_dat = env_array, occ = occ, negative = FALSE)
\end{Sinput}
\begin{Soutput}
[1] -52393.06
\end{Soutput}
\end{Schunk}
\noindent The function \texttt{loglik\_math} first transforms parameters to the 
biological scale using \texttt{math\_to\_bio} and
then evaluates \texttt{loglik\_bio}; so one optimizes it directly 
(see below). If your favorite optimizer minimizes by default, you can use:
\begin{Schunk}
\begin{Sinput}
> loglik_math(param_vector, env_dat = env_array, occ = occ, negative = TRUE)
\end{Sinput}
\begin{Soutput}
[1] 52393.06
\end{Soutput}
\end{Schunk}
\noindent Before moving on to optimizing, we note that \texttt{loglik\_math}
requires a named vector for its input \texttt{param\_vector}. This is to reduce
the possibility of errors stemming form parameter ordering. There is a 
naming convention for arguments which must be followed exactly, and which is described 
in the documentation for \texttt{loglik\_math}. 
The function \texttt{make\_mask\_names} helps the user by giving the parameter names which
are required for a model making use of a given number of environmental variables.

\section{Optimizing the likelihood}

Maximizing the likelihood from one initial parameter guess is straightforward using the built-in
R function \texttt{optim} (and a variety of other optimizers are also available):
\begin{Schunk}
\begin{Sinput}
> optim(par = param_vector,
+       fn = loglik_math,
+       method = "BFGS",
+       env_dat = env_array,
+       occ = occ,
+       negative = TRUE,
+       control = list(trace=100))